problem:  
Let \((M^n, g_0)\) be a complete noncompact Riemannian manifold with Euclidean volume growth and quadratic curvature decay:
\[
\operatorname{Vol}(B_r(p)) \sim C r^n, \qquad |\operatorname{Rm}(g_0)| \leq C(1 + r(p,x))^{-2}.
\]
Assume that the Ricci flow \(\partial_t g = -2\operatorname{Ric}(g)\) with initial data \(g_0\) exists for all \(t > 0\) and satisfies uniform curvature bounds of the form
\[
|\operatorname{Rm}(g(t))| \leq \frac{C}{1+t}.
\]
For any sequence \(t_i \to \infty\), consider the blow-down sequence
\[
(M, g_i(t)) := (M, \lambda_i g(t_i + \lambda_i^{-1}t)), \quad \lambda_i := (1 + t_i)^{-1}.
\]
Under these asymptotic and curvature assumptions, conjecture whether every pointed Cheeger–Gromov limit of \((M, g_i(t), p)\) (for a suitable basepoint \(p\)) is necessarily a gradient expanding Ricci soliton.

Equivalently: **Conjecture** that the combination of Euclidean volume growth, quadratic curvature decay, and Type III curvature bound implies the existence of a canonical expanding soliton model at infinity for noncompact eternal Ricci flow evolutions.

Why is it a "good" problem:  
This problem posits a sharp conjectural link between geometric asymptotics of a noncompact manifold and the self-similar asymptotic behavior of its Ricci flow evolution. It transforms an open-ended program into a concrete, testable assertion—whether Euclidean volume growth and quadratic curvature decay suffice to enforce soliton-like asymptotics for Type III Ricci flows. The question lies precisely at the frontier of singularity analysis on noncompact spaces: current theory describes Type I asymptotics (shrinking solitons) reasonably well, but there is no rigorous understanding of global Type III limits in the noncompact setting. A resolution would unify the geometric large-time structure (asymptotic flatness, curvature decay) with the analytic attractors of the Ricci flow (expanding solitons), potentially revealing a soliton classification parallel to Perelman’s singular model classification. Its formulation is simple yet deeply interconnected with core mechanisms of self-similarity, curvature decay, and long-time flow geometry, making it a concise yet profoundly insightful research problem bridging geometric analysis and global differential geometry.